\documentclass[11pt]{article}

\usepackage[czech]{babel}
\usepackage{a4wide}
\usepackage{fancyhdr}
\usepackage{amssymb}
\usepackage{amsthm}
\usepackage{amsmath}
\usepackage{enumitem}
\usepackage{booktabs}
\usepackage{tabularx}
\usepackage{mathtools}
\usepackage{mleftright}
\usepackage{subcaption}
\usepackage{gensymb}
\usepackage{parskip}
\usepackage{titlesec}
\usepackage[normalem]{ulem}
\usepackage{subfiles}
\usepackage[bottom]{footmisc}  % forcing footnotes to be at the very bottom

\usepackage{geometry}
\geometry{
    a4paper,
    total={170mm,257mm},
    right=20mm,
    left=20mm,
    top=30mm,
    bottom=20mm,
}


% ---------------------------------------- KÓD ----------------------------------------
\usepackage[outputdir=Build]{minted}
\setminted{
framesep=2mm,
linenos,
autogobble,
fontsize=\footnotesize
}

\RecustomVerbatimEnvironment{Verbatim}{BVerbatim}{}
\renewcommand{\figurename}{Listing}


% ---------------------------------------- VLNA ----------------------------------------
\usepackage{luavlna}
\AtBeginDocument{\singlechars{czech}{AaIiVvOoUuSsZzKk}}


% ---------------------------------------- BIBLIOGRAFIE ----------------------------------------
\usepackage[backend=biber, style=numeric]{biblatex}
\addbibresource{main.bib}


% ---------------------------------------- JEDNOTKY ----------------------------------------
\usepackage{siunitx}
\sisetup{per-mode = symbol}


% ---------------------------------------- BOXY ----------------------------------------
\setlength{\fboxrule}{0.5pt}
\setlength{\fboxsep}{0pt}


% ---------------------------------------- FLOATY ----------------------------------------
\usepackage{float}
\floatplacement{figure}{H}


% ---------------------------------------- PŘÍKLADY ----------------------------------------
% https://tex.stackexchange.com/questions/164113/the-exercise-package-in-list-form
\usepackage{scrextend}
\usepackage[load-headings]{exsheets}
\DeclareInstance{exsheets-heading}{mylist}{default}{
  runin = true ,
  attach = {
    main[l,vc]number[l,vc](-2em,0pt) ;
    main[l,vc]title[l,vc](-5.5em,0pt)
  }
}

\DeclareRobustCommand*\questionstar{\texorpdfstring{\bonusquestionsign}{* }}
\DeclareRobustCommand*\bonusquestionsign{\llap{$\bigstar$\space}}

\NewQuSolPair
	{question*}[name=\questionstar Příklad]
	{solution*}[name=\questionstar Řešení]

\DeclareTranslation{Czech}{exsheets-exercise-name}{Příklad}
\DeclareTranslation{Czech}{exsheets-solution-name}{Řešení}

\SetupExSheets{
  headings = mylist , % use the new headings instance
  headings-format = \it ,
  counter-format = se.qu ,
  counter-within = section ,
	question/pre-hook = \addmargin[5.5em]{0em} ,
	solution/pre-hook = \addmargin[5.5em]{0em} ,
	question/post-hook = \endaddmargin ,
	solution/post-hook = \endaddmargin
}

% https://tex.stackexchange.com/questions/263467/exsheets-how-to-get-a-counter-reference-combining-section-and-question-numbers
\renewcommand\thequestion{\thesection.\arabic{question}}


% ---------------------------------------- LINKY ----------------------------------------
\usepackage{hyperref}
\urlstyle{same}
\hypersetup{
    colorlinks=true,
    urlcolor=blue,
    linkcolor=black,%TOC
}


% ---------------------------------------- NADPISY ----------------------------------------
\titleformat{\section} {\normalfont\fontsize{16}{15}\bfseries}{\thesection}{1em}{}
\titleformat{\subsection} {\normalfont\fontsize{14}{15}\bfseries}{\thesubsection}{1em}{}
\titleformat{\subsubsection} {\normalfont\fontsize{12}{15}\bfseries}{\thesubsubsection}{1em}{}

\usepackage{tcolorbox}
\tcbset{on line,
        boxsep=4pt, left=0pt,right=0pt,top=0pt,bottom=0pt,
        colframe=white
        }


% ---------------------------------------- HLAVIČKA ----------------------------------------
\pagestyle{fancy}
\fancyhf{}
% Now set the headers
\fancyhead[LE,RO]{\footnotesize\thepage}
\fancyhead[LO,RE]{\footnotesize\leftmark}
\fancyfoot{}
\fancyfoot[R]{\thepage}

% ---------------------------------------- BARVY ----------------------------------------
\definecolor{ceruleanblue}{rgb}{0.16, 0.32, 0.75}
\definecolor{motorlight}{rgb}{0.85, 0.80, 0.70}
\definecolor{motordark}{rgb}{0.64, 0.51, 0.25}
\definecolor{wait}{rgb}{0.64, 0.25, 0.25}


% ---------------------------------------- DEFINICE BLOKŮ ----------------------------------------
\newcommand{\block}{$\vcenter{\hbox{\includegraphics[height=0.8em]{Images/blocky-logo.png}}}$}
\newcommand{\centerimage}[2]{$\vcenter{\hbox{\includegraphics[height=#1]{#2}}}$}

% TODO: sizes to variable
% TODO: do this more programatically...
\newcommand{\blockMotorStart}{\item[\block] \centerimage{2em}{Images/blocks/motor-start.png} -- \parbox{0.545\textwidth}{zapne daný motor dopředu (\textit{forward}) nebo dozadu (\textit{reverse}) na danou rychlost (\textit{velocity}) v procentech.}}
\newcommand{\blockMotorStop}{\item[\block] \centerimage{1.6em}{Images/blocks/motor-stop.png} -- vypne daný motor.}
\newcommand{\blockMotorDistanceImage}{\centerimage{2em}{Images/blocks/motor-distance-start.png} }
\newcommand{\blockMotorDistance}{\item[\block] \blockMotorDistanceImage -- \parbox{0.525\textwidth}{zapne motor, dokud se neotočí o daný počet stupňů (\textit{degrees}) nebo otáček (\textit{revolutions}). Je neblokující -- zapne motor a jde na další block.}}
\newcommand{\blockMotorVelocityImage}{\centerimage{2em}{Images/blocks/motor-velocity.png} }
\newcommand{\blockMotorVelocity}{\item[\block] \blockMotorVelocityImage -- nastaví rychlost motoru, ale nezapne ho!}

\newcommand{\blockMotorDoneImage}{\centerimage{1.6em}{Images/blocks/motor-done.png} }
\newcommand{\blockMotorDone}{\item[\block] \blockMotorDoneImage -- podmínka která říká, zda daný motor skončil.}

\newcommand{\blockLoop}{\item[\block] \centerimage{4em}{Images/blocks/loop.png} -- opakuje bloky uvnitř několikrát.}
\newcommand{\blockLoopForever}{\item[\block] \centerimage{4em}{Images/blocks/loop-forever.png} -- opakuje bloky uvnitř donekonečna.}
\newcommand{\blockLoopWhile}{\item[\block] \centerimage{4em}{Images/blocks/loop-while.png} -- \parbox{0.725\textwidth}{opakuje bloky uvnitř, dokud (\textit{while}) nebo než (\textit{until}) platí podmínka.}}

\newcommand{\blockWaitImage}{\centerimage{2em}{Images/blocks/wait.png} }
\newcommand{\blockWait}{\item[\block] \blockWaitImage -- čeká daný čas.}
\newcommand{\blockWaitUntil}{\item[\block] \centerimage{1.6em}{Images/blocks/wait-until.png} -- čeká, než (\textit{until}) platí podmínka.}
\newcommand{\blockStartImage}{\centerimage{1.6em}{Images/blocks/start.png} }
\newcommand{\blockStart}{\item[\block] \blockStartImage -- startovní blok pro každý Blocky program.}
\newcommand{\blockComment}{\item[\block] \centerimage{1.6em}{Images/blocks/comment.png} -- blok, který nic nedělá (pouze komentuje kód).}

\newcommand{\blockFunctionDefinition}{\item[\block] \centerimage{4em}{Images/blocks/function-definition.png} -- definování funkce; jméno lze upravit kliknutím.}
\newcommand{\blockFunctionCall}{\item[\block] \centerimage{2em}{Images/blocks/function-call.png} -- volání funkce TODO lepší vysvětlení?.}

\newcommand{\blockBumperPressed}{\item[\block] \centerimage{1.6em}{Images/blocks/bumper-pressed.png} -- podmínka, která říká, zda je daný senzor stisknutý.}

\newcommand{\blockVariableCreate}{\item[\block] \centerimage{1.6em}{Images/blocks/variable-create.png} -- vytvoří novou proměnnou.}
\newcommand{\blockVariableChange}{\item[\block] \centerimage{1.6em}{Images/blocks/variable-change.png} -- změní danou proměnnou o dané číslo.}
\newcommand{\blockVariableGet}{\item[\block] \centerimage{1.6em}{Images/blocks/variable-get.png} -- hodnota dané proměnné.}
\newcommand{\blockVariableSet}{\item[\block] \centerimage{1.6em}{Images/blocks/variable-set.png} -- nastaví hodnotu dané proměnné.}

\newcommand{\blockMathOperation}{\item[\block] \centerimage{2em}{Images/blocks/math-operation.png} -- spočítá výsledek dané matematické operace.}
\newcommand{\blockMathTest}{\item[\block] \centerimage{2em}{Images/blocks/math-test.png} -- podmínka pro běžné matematické testy (sudé/liché, kladné/záporné...). .}
\newcommand{\blockMathValue}{\item[\block] \centerimage{1.6em}{Images/blocks/math-value.png} -- číselná hodnota.}
\newcommand{\blockMathConstant}{\item[\block] \centerimage{1.6em}{Images/blocks/math-constant.png} -- hodnota běžných matematických konstant ($\pi, e, \sqrt{2}\ldots$).}
\newcommand{\blockMathRandom}{\item[\block] \centerimage{2em}{Images/blocks/math-random.png} -- hodnota náhodného čísla v daném rozmezí.}

\newcommand{\blockLogicTrueImage}{\centerimage{1.6em}{Images/blocks/logic-true.png}}
\newcommand{\blockLogicTrueImageBaseline}{\centerimage{\baselineskip}{Images/blocks/logic-true.png}}
\newcommand{\blockLogicTrue}{\item[\block] \blockLogicTrueImage\ -- hodnota pravdy.}
\newcommand{\blockLogicFalseImage}{\centerimage{1.6em}{Images/blocks/logic-false.png}}
\newcommand{\blockLogicFalseImageBaseline}{\centerimage{\baselineskip}{Images/blocks/logic-false.png}}
\newcommand{\blockLogicFalse}{\item[\block] \blockLogicFalseImage\ -- hodnota neprady.}

\newcommand{\blockString}{\item[\block] \centerimage{1.6em}{Images/blocks/string.png} -- text.}
\newcommand{\blockDisplayPrint}{\item[\block] \centerimage{2em}{Images/blocks/display-print.png} -- vypsání textu na daný řádek.}
\newcommand{\blockDisplayClear}{\item[\block] \centerimage{2em}{Images/blocks/display-clear.png} -- smazání textu na daném řádku.}


% ---------------------------------------- SDÍLENÉ PŘÍKAZY ----------------------------------------
\newcommand{\errata}{%
	\subsection{Chyby, překlepy, doplňky}
	Pokud někde odhalíte chybu/překlep, nebo máte pocit, že by něco šlo napsat lépe, tak mě buďto kontaktuje na adrese \href{mailto:tomas@slama.dev}{tomas@slama.dev}, nebo vytvořte issue/pull request na stránce \url{https://github.com/xiaoxiae/uvod-do-programovani-vex-v5}.
}


\begin{document}
\begin{titlepage}
    \begin{center}
        \vspace*{3em}
        \Huge
				\textbf{Metodika k programování} \includegraphics[height=0.65\baselineskip]{Images/vex-logo.png}
				\\
				\Large
        \vspace*{0.5em}
				v jazyce Blocky a prostředí RobotMesh studio

        \vfill

        \flushright
        \normalsize
				Tomáš Sláma\\
				\today
    \end{center}
\end{titlepage}

\tableofcontents
\clearpage

\setcounter{secnumdepth}{0}

\section{Na úvod}
Tento krátký dokument je doprovodný text pro učebnici \textbf{Úvod do programování \includegraphics[height=0.65\baselineskip]{Images/vex-logo.png}.} Jejím cílem je poskytnout dodatečné informace vyučujícímu, který používá učebnici jako hlavní zdroj informací a příkladů.

\errata

\newpage

\section{Příprava}
V rámci přípravy na každou hodinu je důležité si prostudovat materiál dané kapitoly a ujistit se, že bude jeho podání pro studenty dostatečně pochopitelné a srozumitelné.

Příprava je obzvlášť důležitá u příkladů -- kromě pochopení řešení je nutné si připravit věci na jejich úspěšné řešení (papírové ruličky/špejle u robotího slalomu, papíry/náhradní díly robota při vytváření bludiště, apod.)

\section{Příklady a řešení}
Řešení příkladů je hlavní mechanizmus, kterým se studenti naučí a opravdu pochopí probíranou látku. Při řešení je však nesmírně důležité, aby na řešení přišli \textbf{sami} -- vyzrazením řešení se student prakticky nic nenaučí. Při zaseknutí je výrazně efektivnější studenta popostrčit správným směrem, než mu vyzradit řešení:
\begin{itemize}
	\item chápeš, na co přesně se tě příklad ptá?
	\item zkoušel jsi použít tento blok?
	\item jsi si jistý, že tato část programu dělá to, co si myslíš?
\end{itemize}

\vspace{\baselineskip}

Pro studenty, kteří jsou znatelně rychlejší než ostatní, je dobré vymýšlet obtížnější variace na problémy, které aktuálně řeší:
\begin{itemize}
	\item robot měl jezdit do čtverce -- ať jezdí do trojúhelníku
	\item robot se měl otočit doprava -- ať se otočí doprava s použitím jednoho kola
	\item robot se měl otočit přesně o 90\degree -- ať se otočí o 45\degree
\end{itemize}

\section{Způsob výkladu}
V průběhu hodiny budou studenti mít k dispozici notebooky/počítače, na kterých budou robota programovat. U samotného výkladu (který bude ukazován na projektoru/monitoru) je ale nejlepší, když je nepoužívají a plně se soustředí na vás -- látku si zapamatují lépe, poslouchat a zkoušet to co vyučující říká zároveň nejde.

Tam, kde je to možné, je nejlepší látku ukazovat přímo na příkladech na samotném robotovi. Vědět to, co by program teoreticky měl dělat vždy je výrazně slabší než program vidět v akci.

Atmosféru hodiny ovlivňuje především hlas a výraz učitele -- energický výklad z úsměvem pomáhá přenést zapálení i na ostatní, stejně jako monotónní výklad se znuděným výrazem má tendenci uspávat.

\section{Průběh běžné hodiny}
V rámci struktury hodiny je dobré se držet stejného formátu.

Rozhodně není nutné dodržovat formát rozepsaný níže -- je skvělé experimentovat (v rozumné míře, ať z hodiny studenti něco mají) s různými formáty a najít ten, který vám osobně sedí nejvíc. Mě se osobně osvědčilo následující:

\begin{table}[h!]
	\begin{tabularx}{\textwidth}{rc>{\setlength{\parskip}{0.5\baselineskip}}X} \toprule
		\textbf{úvod} & $2$ min. & Krátká otázka o tom, co všichni přes týden dělali/co je zaujalo/nad čím přemýšleli. Pomáhá studenty „probudit“ a připravit na to se něco nového naučit. Krátká diskuze/debata je naprosto v pořádku.

		Na první hodině je úvod spíše brát jako seznámení -- učitel začíná, po kroužku všichni postupně zodpoví (klidně s kolujícím předmětem):
		\begin{itemize}[label=\textbullet, leftmargin=20pt,
							after*={\mbox{}\vspace{-\baselineskip}}]
			\item jméno/přezdívka
			\item oblíbený kroužek/aktivita
			\item co od kurzu očekávají (včetně učitele!)
		\end{itemize}
		\\
		\noalign{\vspace{8pt}}
		\textbf{opakování} & $5$ min. & Zopakování toho, co se dělalo minulou hodinu (hlavně obtížnějších konceptů). Prostor na pokládání otázek. Na první hodině nedává smysl, je dobré tuto část přeskočit. \\
		\noalign{\vspace{8pt}}
		\textbf{výuka} & zbytek &  Hlavní část hodiny. Probírání nové látky, plnění příkladů. \\
		\noalign{\vspace{8pt}}
		\textbf{zakončení} & $5$ min. & Rychlé shrnutí toho, co se za dnešek dělalo. Prostor na dotazy k dnešní látce a drobné nastínění toho, co se bude dělat příští hodinu. \\
		\bottomrule
	\end{tabularx}
\end{table}

Není nutné dodržovat stejný formát na všech hodinách, protože dost záleží na jejich obsahu (pokud celou hodinu probíhala soutěž o programování pohybu robota v bludišti, tak její opakování je poněkud zbytečné).

\end{document}
